% answers_short.tex —— 2010 年·普通高等学校招生考试·北京卷·理科 · 简短答案速查
% 由 tools/make_short_answers.py 自 answers.tex 的 \gksolhead/\gkanswer 提取，
% 并与 2010/2010_数学真题_北京理_答案要点.md 逐题比对后生成，勿手改（改源 MD 后重跑）。

\begin{gknotice}
\textbf{使用说明}：本节只给最终答案便于快速核对，完整推导见解析版。答案为编写者独立推导，\textbf{非官方标准答案}。
\end{gknotice}

\gkpart{选择题}
\begin{center}\begin{tabular}{ccccccccc}
\toprule
题号 & 1 & 2 & 3 & 4 & 5 & 6 & 7 & 8 \\
\midrule
答案 & B & C & B & A & C & B & A & D \\
\bottomrule
\end{tabular}\end{center}

\gkpart{填空题}
\begin{itemize}[leftmargin=2.6em]
  \item[\textbf{9.}] $(-1,1)$。
  \item[\textbf{10.}] $1$。
  \item[\textbf{11.}] $0.030$；$3$。
  \item[\textbf{12.}] $5$；$2\sqrt7$。
  \item[\textbf{13.}] $(\pm4,0)$；$y=\pm\sqrt3x$。
  \item[\textbf{14.}] $4$；$\dfrac\pi2+3$。
\end{itemize}

\gkpart{解答题}
\begin{itemize}[leftmargin=2.6em]
  \item[\textbf{15.}] （1）$-\dfrac14$；（2）最大值 $2$，最小值 $-1$。
  \item[\textbf{16.}] （1）（2）证明；（3）$\dfrac\pi6$。
  \item[\textbf{17.}] （1）$\dfrac{119}{125}$；（2）$p=\dfrac35$，$q=\dfrac25$；（3）$\dfrac95$。
  \item[\textbf{18.}] （1）$y-\ln2=\dfrac32(x-1)$；（2）单调区间见解析。
  \item[\textbf{19.}] （1）$\dfrac{x^2}4+\dfrac{3y^2}4=1$（$x\neq\pm1$）；（2）$\left(\dfrac53,\pm\dfrac{\sqrt{33}}9\right)$。
  \item[\textbf{20.}] 三问均证明。
\end{itemize}
